\documentclass[11pt,a4paper]{article}

% --- PACKAGES ---
\usepackage[utf8]{inputenc}
\usepackage[english]{babel}
\usepackage{geometry}
\geometry{margin=2.5cm}
\usepackage{graphicx} % Per le immagini
\usepackage{listings} % Per il codice C++
\usepackage{xcolor}   % Per i colori del codice
\usepackage{hyperref} % Per i link cliccabili
\usepackage{amsmath}  % Per le formule matematiche

% --- CODE STYLE CONFIGURATION ---
\definecolor{codegreen}{rgb}{0,0.6,0}
\definecolor{codegray}{rgb}{0.5,0.5,0.5}
\definecolor{codepurple}{rgb}{0.58,0,0.82}
\definecolor{backcolour}{rgb}{0.95,0.95,0.92}

\lstset{
    backgroundcolor=\color{backcolour},   
    commentstyle=\color{codegreen},
    keywordstyle=\color{magenta},
    numberstyle=\tiny\color{codegray},
    stringstyle=\color{codepurple},
    basicstyle=\ttfamily\footnotesize,
    breaklines=true,
    captionpos=b,                    
    keepspaces=true,                 
    numbers=left,                    
    showspaces=false,                
    showstringspaces=false,
    showtabs=false,                  
    tabsize=2,
    language=C++
}

% --- DOCUMENT INFO ---
\title{\textbf{Technical Report: Smart-Hysteresis Thermal Monitor}}
\author{Alessio Nuzzolese}
\date{}

\begin{document}

\maketitle

\section{Introduction}
This report documents the development of a temperature control system based on the \textbf{Arduino Uno} platform. The system is designed to provide real-time thermal monitoring with an integrated safety alarm governed by hysteresis logic to ensure signal stability.

\section{System Architecture}
The project was prototyped using the \textit{Autodesk Tinkercad} environment. The hardware components, all included in the \textbf{Elegoo Super Starter Kit}, are:
\begin{itemize}
    \item \textbf{Microcontroller:} Arduino Uno R3.
    \item \textbf{Sensing:} TMP36 Precision Analog Temperature Sensor.
    \item \textbf{Display:} I2C 16x2 Liquid Crystal Display (LCD).
    \item \textbf{Indicators:} Red LED and Passive Piezoelectric Buzzer.
\end{itemize}

\begin{figure}[h]
    \centering
  
    \includegraphics[width=0.8\textwidth]{Thermometer.png} 
    \caption{TinkerCad circuit view.}
    \label{fig:circuit}
\end{figure}

\begin{figure}[h]
    \centering
  
    \includegraphics[width=0.8\textwidth]{Thermometer2.png} 
    \caption{TinkerCad wiring schematic view.}
    \label{fig:circuit}
\end{figure}


\section{Mathematical Modeling}
\subsection{Analog-to-Digital Conversion}
The Arduino ADC converts the analog voltage from the TMP36 sensor into a 10-bit digital value ($0-1023$). The conversion to Celsius degrees follows the transfer function:
\begin{equation}
    T (^\circ C) = \left( \frac{ADC_{raw} \cdot 5.0}{1024} - 0.5 \right) \cdot 100
\end{equation}
This formula accounts for the $500mV$ offset at $0^\circ C$ and the $10mV/^\circ C$ scale factor typical of the TMP36 sensor.

\subsection{Hysteresis Logic}
To prevent "chattering" (rapid switching of the alarm due to noise), a \textbf{Schmitt Trigger} approach was implemented with two distinct thresholds:
\begin{itemize}
    \item \textbf{Activation ($S_{on}$):} $30.0^\circ C$
    \item \textbf{Deactivation ($S_{off}$):} $28.0^\circ C$
\end{itemize}

\section{Implementation}
The control firmware was written in C++. Below is the implementation of the hysteresis logic:

\begin{lstlisting}
// --- HYSTERESIS LOGIC ---
if (temperatureC >= activation) {
    alarmOn = true;  // Trigger alarm
} else if (temperatureC <= deactivation) {
    alarmOn = false; // Reset alarm
}
\end{lstlisting}

\section{Conclusion}
The simulation successfully demonstrated that the hysteresis logic provides a stable control signal, effectively filtering out minor temperature fluctuations. This project serves as a foundational example of industrial sensor interfacing and robust software design.
This project has real application in every engineering field.

\end{document}